% !TEX root = ../main.tex
% !TEX program = pdflatex
\subsection{Определение размера каждого уровня кэша}

\textbf{Метод.} Для оценки L1 строится случайная цепочка с одним узлом на каждые 64~байта; для отдельной оценки L2 запускается второй sweep с одним узлом на каждые 128~байт. Обе конфигурации исследуются на диапазоне от 4~КБ до 256~МБ. В каждом sweep измеряется минимум, среднее и медиана пяти замеров с двумя прогревочными проходами перед каждым замером. Длина измеряемого прохода охватывает всю цепочку не менее одного раза.

Шаг 64~байта соответствует размеру строки L1, сообщаемому экспериментом. Шаг 128~байт позволяет во втором sweep посещать по одному узлу на заявленный ОС размер строки. Это снижает неоднозначность фактически затронутого объёма, но не устраняет влияние трансляции виртуальных адресов и конфликтов кэша.

\textbf{Результаты} показаны на рисунке~\ref{fig:size} и в таблице~\ref{tab:size}. На графике приведены минимум и среднее; медиана сохранена в CSV для повторной обработки.

\fig{size.png}{Время доступа для цепочек с шагом 64 и 128 байт}{fig:size}

\begin{table}[H]
\centering
\caption{Медианное время доступа в размерных sweep, нс}\label{tab:size}
\begin{tabular}{lccc}
\toprule
Адресный диапазон & Шаг 64 Б & Шаг 128 Б & Наблюдение \\
\midrule
128 КБ & 1,15 & 1,15 & задержка L1 \\
160 КБ & 3,49 & 4,29 & переход от L1 \\
224 КБ & 5,42 & 5,62 & задержка L2 \\
4 МБ & 5,59 & 5,89 & плато L2 \\
8 МБ & 6,36 & 6,57 & начало роста \\
16 МБ & 12,24 & 16,30 & смешанная область \\
24 МБ & 76,00 & 53,02 & существенные промахи \\
256 МБ & 102,21 & 106,27 & основная память \\
\bottomrule
\end{tabular}
\end{table}

\textbf{Анализ.}
\begin{itemize}
  \item В sweep с шагом 64~байта задержка остаётся около 1,14~нс до 128~КБ и возрастает при 160~КБ. Это согласуется с размером L1 данных P-ядер 128~КБ, указанным ОС.
  \item Для шага 128~байт плато около 5,6--5,9~нс видно от 224~КБ до 4~МБ. Оно соответствует попаданиям в L2. Рост задержки начинается примерно с 8~МБ и становится резким между 16 и 24~МБ.
  \item ОС сообщает размер L2 P-кластера 16~МБ. Размерный эксперимент согласуется с кэшем такого масштаба, однако не показывает точную границу одним скачком: результат зависит также от TLB, физического распределения страниц и общего кэша SLC. Поэтому 16~МБ приводится как значение ОС, а кривая — как экспериментальная проверка порядка величины, а не самостоятельное точное измерение ёмкости.
\end{itemize}

На больших наборах задержка составляет 102--106~нс. Плавный переход объясняют промахи TLB и попадания в SLC.

Класс обслуживания \code{QOS_CLASS_USER_INTERACTIVE} не закрепляет поток за P-ядром, так что отдельные точки могут отражать работу на разных типах ядер.
